\section{CompCert Case Study: Optimistic Affine Loops}
\label{sec:case-study}

The case study tests the interfaces against real Clight execution and memory.
PolCert supplies a capability reference for affine domains, candidate schedules,
and tiling; the integration uses GuardCert's own source/guard adapters together
with the existing vendored optimizer proof infrastructure. The concrete toolchain
is pinned to CompCert~3.18, Rocq~9.2, and Stdlib~9.2.

\subsection{Established Compiler Path}

The current complete path supports a loaded signed root bound with a constant
offset and canonical affine descendants. It captures the initial header word,
executes private numeric and memory checks, validates an untrusted candidate,
and retains the original source on refusal. Candidates include schedule
reordering and supported tiling/domain witnesses; successful checking binds the
candidate to the actual source description.

For this path, \code{compile_offset_affine_multi_regions_correct} proves that a
successful compiler result is related to the original Csem program by a
backward simulation into CompCert assembly semantics. The proof composes
SimplExpr, SimplLocals, the expression-region host, and the verified backend.
The region contracts and host progress obligations determine the source class;
the theorem is not a license to install arbitrary Clight fragments.

\subsection{A Nested Loaded-Bound Acceptance Target}

An adapted Figure~2 source from OLO~\cite{olo2017} supplies a concrete pending
acceptance target:
\begin{lstlisting}[language=C]
for (row = 0; row < shape[0] + 1; row++)
  for (column = 0; column < shape[1] + 1; column++)
    for (component = 0; component < 5; component++)
      output[(row * 16 + column) * 5 + component] =
          row + column + component;
\end{lstlisting}
This preserves the three-level control pattern and two repeatedly loaded
headers. It uses a flat fixed layout and an arithmetic store body; it is not
the full NAS BT kernel. The output may alias either header. A store can then
change the next source test, so the captured iteration domain need not describe
the original execution until preservation has been established.

The language services first capture the outer header. They capture the child
header only if the actual outer test is active. An empty outer loop therefore
requires neither a defined child counter nor a readable second header. The
source's modular \code{Int.add} computation is kept distinct from the raw memory
observation; no-wrap facts belong to the accepted modeling condition.

For an active prefix, the proof retains a witness in the memory actually reached
by the source. It transports access permissions back to the guard-entry memory,
not loaded data values. Each accepted body check must protect \emph{all} captured
observations before the next child test is justified. Completion of those checks
for the current row permits the next outer test. Only after this preservation
chain can the proof derive execution of the cached source and consume the
affine candidate certificate.

The ordered capture, indexed-load primitives, nested source-prefix services,
and two-cache transport are proved. A numeric probe can already run from
captured parameter words without assuming completion of the cached source.
Additional adapters now initialize private model bounds and transport execution
of a reached literal-bound subbody into the canonical affine model. They derive
that subdomain's access permissions at guard entry while retaining the reached
source memory. This bridge requires no stability assumption for the enclosing
loaded headers.

The domain adapter now consumes these permissions to execute a joint scan for
the reached constant-bound subbody. Observation descriptors carry actual address
expressions and load receipts, permitting comparisons against both
\code{shape} and \code{shape + 1} without another pointer temporary. Acceptance
establishes separation of every point write from every captured word, then
preservation for actual subbody executions with the same coordinate and pointer
view. This evidence discharges the inner-prefix preservation obligation.
A concrete fixture proves all five original stores and the actual scan: data
and headers overlapping are refused, while adjacent disjoint regions within one
allocation are accepted. It does not install an enclosing loop optimizer.

The subsequent adapter executes the inner short-circuit loop using source
prefix receipts for each column. Its acceptance supplies observation preservation
for every column and advances the current outer prefix. The source witness and
public guard state remain separate; mapped private cursors provide the scanned
coordinates. A refusal at the first body leaves the scan cursor at zero, and an
empty-child execution needs no body or output-pointer receipt. Fixed-entry prefix
services require the header law only at the captured entry. Numeric-domain,
word-view, and namespace facts remain explicit inputs until the factory produces
them from the original source and checked package.

The outer adapter now consumes this row producer in an actual short-circuit
loop. Complete acceptance supplies observation preservation for every active
row and column, discharging the preservation input to the existing two-cache
transport. The resulting cached source has the same exit temporaries and final
memory as the original source. A first-row refusal stops before the outer
increment. Separate empty-outer execution and zero-header transport services
need no child cache or body-pointer receipt. An unread child must not become a
premise merely to instantiate a two-observation prefix. These services do not
by themselves license a reordered candidate that reads the child first.
\pending{Assemble the canonical three-level affine model, the actual
capture/numeric producers for the typed metadata inputs, the original-source matcher/factory, and a new
whole-program compiler entry. The current adapted Figure~2 native probe retains
the original source and does not demonstrate an installed optimization.}

\subsection{Condition Expressiveness and Reuse}

The supported services cover machine-word ranges/no-wrap, byte-level memory
separation, permissions for typed accesses, and preservation of loaded values.
These are sufficient-condition libraries with declared syntax and semantic
domains. They do not decide arbitrary propositions or infer weakest conditions
for arbitrary program pairs. For a source-derived scan, the domain implementation
still proves access coverage and the source/model bridge. The language library
supplies the physical comparison and load/store preservation lemmas.

Condition construction must preserve partial evaluation safety. A mathematically
simpler condition can read a header or compare a pointer earlier than the source
permits. Reuse therefore requires a certificate for the reordered or simplified
check, not just Boolean equivalence on states where every expression is defined.
\pending{Select the final compact stability condition and compare its guard
work and useful acceptance against the existing scan on the same source. General
projection and arbitrary dependent preload expressions remain outside the
established implementation claims.}
